\makeatletter
\def\input@path{{./}{../}{../../}{../../../}{../../../../}{preamble/}{../preamble/}{../../preamble/}{../../../preamble/}{../../../../book/preamble/}}
\makeatother
\input{preamble_exams.tex}
\input{graphs/highlights.tex}
% \GTenHighlightsfalse % Uncomment to disable instructional highlights.
% Population/model, sample evidence, design threshold, and result roles.
\newcommand{\Pop}[1]{\HighlightStroke{GTenHLParameter}{#1}}
\newcommand{\Obs}[1]{\HighlightStroke{GTenHLValue}{#1}}
\newcommand{\Design}[1]{\HighlightStroke{GTenHLThird}{#1}}
\newcommand{\Result}[1]{\HighlightStroke{GTenHLOutcome}{#1}}

\title{Grade 11 Global Economics\\Term 3 --- Lesson 1\\Home Practice C7: Paired-sample t-tests}
\author{}\date{}
\hypersetup{pdftitle={Grade 11 Global Economics - Term 3 - Lesson 1 - C7},pdfsubject={Paired-sample t-tests}}
\begin{document}
\maketitle
\setcounter{tocdepth}{1}
\setcounter{secnumdepth}{2}
\phantomsection\label{toc}
\tableofcontents
\section*{Evaluated criteria}
\begin{description}
\item[\textbf{C7}] Conducts paired-sample t-tests to evaluate mean differences between matched or repeated observations, verifying the required conditions and interpreting the results in economic and financial contexts.
\end{description}
\textbf{How to practise.} Write the difference direction first. Work with one difference per pair and count pairs, not the total measurements. Attempt each task before opening its solution link. Use the return arrows to revisit the problem, group, or contents. All datasets and scenarios in this activity are illustrative.
\clearpage
\section{Introduction}\label{sec:introduction}
\subsection*{Highlight legend}
\phantomsection\label{tab:highlight-legend}
\begin{tabularx}{\linewidth}{@{}>{\raggedright\arraybackslash}p{.22\linewidth}>{\raggedright\arraybackslash}p{.25\linewidth}>{\raggedright\arraybackslash}X@{}}
\toprule
\textbf{Symbol or notation}&\textbf{Concept}&\textbf{Brief meaning}\\\midrule
$\Obs{d_i}$ & \Obs{Paired difference} & One value per matched pair. \\
$\Pop{\mu_d},\Obs{\bar d}$ & \Pop{Mean difference} & Population effect and estimate. \\
$\Obs{s_d},\Obs{n}$ & \Obs{Spread and pair count} & Variability of differences. \\
$\Obs{t},\Design{df}$ & \Design{Paired t model} & Statistic and reference. \\
\bottomrule
\end{tabularx}
\subsection{Introductory example 1: Two observations on one shop}\label{ex:intro-one}
A shop's weekly transfer costs are 30 USD before a new account and 28 USD after. Define $\Obs{d_i}=\text{before}-\text{after}=2$ USD. The observations belong together; two independent groups would discard this match.
\subsection{Introductory example 2: Nine matched shops}\label{ex:intro-two}
Nine shops supply nine complete pairs with $\Obs{\bar d}=2$ USD and $\Obs{s_d}=3$ USD. The population question is whether $\Pop{\mu_d}>0$, not whether either before or after mean alone is zero. A paired t model uses eight degrees of freedom.
\section{Notation Reference and Practice Procedure}\label{sec:notation-and-procedure}
\subsection{Notation reference}\label{tab:notation-reference}
\begin{tabularx}{\linewidth}{@{}>{\raggedright\arraybackslash}p{.16\linewidth}>{\raggedright\arraybackslash}p{.18\linewidth}>{\raggedright\arraybackslash}p{.28\linewidth}>{\raggedright\arraybackslash}X@{}}
\toprule
\textbf{Symbol or notation}&\textbf{Concept}&\textbf{Definition}&\textbf{Meaning in the introductory examples}\\\midrule
$\Obs{d_i}$ & \Obs{Paired difference} & Here $d_i=\text{before}_i-\text{after}_i$ for cost; positive means a saving. & \textbullet\ \hyperref[ex:intro-one]{Example 1}: $30-28=2$ USD. \\
$\Pop{\mu_d},\Obs{\bar d}$ & \Pop{Mean difference} & $\mu_d$ is the population mean of differences; $\bar d=\sum d_i/n$ estimates it. & \textbullet\ \hyperref[ex:intro-one]{Example 1}: population mean cost saving.\par\textbullet\ \hyperref[ex:intro-two]{Example 2}: sample mean saving 2 USD. \\
$\Obs{s_d},\Obs{n}$ & \Obs{Spread and pair count} & $s_d=\sqrt{\sum(d_i-\bar d)^2/(n-1)}$; $n$ is the number of complete pairs. & \textbullet\ \hyperref[ex:intro-two]{Example 2}: $s_d=3$ USD and $n=9$ pairs. \\
$\Obs{t},\Design{df}$ & \Design{Paired t model} & $t=(\bar d-\mu_{d,0})/(s_d/\sqrt n)$ with $df=n-1$. & \textbullet\ \hyperref[ex:intro-two]{Example 2}: $df=8$; the benchmark is zero saving. \\
\bottomrule
\end{tabularx}
\subsection{Practice procedure}
\begin{enumerate}
\item Define the population of pairs and the order of subtraction. State the benchmark mean difference and alternative before inspecting results.
\item Match observations correctly; each unit contributes one difference. Check random sampling of pairs and independence between pairs.
\item Assess the distribution of differences. Small $n$ requires approximately normal differences without severe outliers; individual before and after distributions are not the relevant normality condition.
\item Compute $\bar d$, $s_d$, $SE=s_d/\sqrt n$, $t$, and $df=n-1$. Use the planned tail, critical value, and p-value.
\item State the decision and translate the sign into the context. Separate evidence about change from a causal claim.
\end{enumerate}
\section{General Worked Problems}\label{sec:general-worked}
\SectionProblemsTableOfContents{\SectionProblemLink{prob:raw}{Five shops and account costs} \\
\SectionProblemLink{prob:direction}{Reversing the subtraction}}
\subsection{Five shops and account costs}\label{prob:raw}
Five randomly selected independent shops record weekly transfer costs in USD:
\begin{center}\begin{tabular}{lrrrrr}\toprule Shop&1&2&3&4&5\\\midrule Before&30&40&50&60&70\\After&28&36&44&52&60\\\bottomrule\end{tabular}\end{center}
Assume the population of cost differences is approximately normal without serious outliers. Test a reduction at $\alpha=0.05$. Supplied for $df=4$: right critical 2.132 and $P(T_4\ge4.243)\approx0.0066$.
\QuestionsTasks[sec:general-worked]{{7}/{sol:raw}/{Construct the differences and conduct the paired test.}}
\subsubsection*{Worked solution}\label{sol:raw}
Use $d=\text{before}-\text{after}$, so a reduction means positive saving. Parameter $\Pop{\mu_d}$ is population mean weekly saving; $H_0:\mu_d=0$, $H_a:\mu_d>0$. The sample consists of five independent pairs; normality is required of their differences and is assumed. Differences are $2,4,6,8,10$ USD:
\[
\bar d=\frac{30}{5}=6,\quad
s_d^2=\frac{(-4)^2+(-2)^2+0^2+2^2+4^2}{4}=10,
\]
\[
s_d=\sqrt{10},\quad SE=\sqrt{10}/\sqrt5=\sqrt2,\quad
t=\frac{\Obs{6}-\Pop{0}}{\sqrt2}\approx4.243,\quad df=4.
\]
Since $4.243>2.132$ and p-value $0.0066<0.05$, \Result{reject $H_0$}. Evidence supports positive population mean savings, estimated as 6 USD weekly per shop. Ten individual measurements still give only five pairs. Without randomized treatment and control, time changes could also explain a before--after difference.
\secInicio[prob:raw][sec:general-worked]
\subsection{Reversing the subtraction}\label{prob:direction}
For nine randomly sampled independent shops, define before-minus-after cost differences. The summary is $\bar d=2$ USD, $s_d=3$ USD, with approximately normal differences and no serious outliers. Test a reduction at $\alpha=0.05$. Supplied for $df=8$: critical values $\pm1.860$ for one-tail 5\%, and $P(T_8\ge2)=P(T_8\le-2)\approx0.0403$.
\QuestionsTasks[sec:general-worked]{{7}/{sol:direction}/{Test with the stated direction, then reverse every difference and explain what changes.}}
\subsubsection*{Worked solution}\label{sol:direction}
With before-minus-after, $H_0:\mu_d=0$, $H_a:\mu_d>0$. Conditions apply to the nine independent differences. $SE=3/\sqrt9=1$, so $t=(2-0)/1=2$, $df=8$. Since $2>1.860$ and p-value $0.0403<0.05$, \Result{reject $H_0$}; evidence supports mean cost reduction. If $e=\text{after}-\text{before}$, then $\bar e=-2$, $s_e=3$, $H_a:\mu_e<0$, and $t=-2$. Compare with $-1.860$ and use the same p-value. The conclusion remains unchanged. Reversing signs without reversing the alternative would ask a different economic question.
\secInicio[prob:direction][sec:general-worked]
\section{Advanced Problems}\label{sec:advanced}
\SectionProblemsTableOfContents{\SectionProblemLink{prob:conditions}{Missing pairs and an influential difference}}
\subsection{Missing pairs and an influential difference}\label{prob:conditions}
A bank measures processing time before and after software installation for the same 12 branches. Three after-measurements are missing because those branches had severe delays. Among the nine complete branches, one difference is much larger than the others; the remaining differences are clustered near zero. The manager wants to treat the 12 before and nine after values as independent samples.
\QuestionsTasks[sec:advanced]{{7}/{sol:conditions}/{Choose the data structure, explain the conditions at risk, and propose a justified next step without inventing a test result.}}
\subsubsection*{Worked solution}\label{sol:conditions}
The unit is a branch measured twice, so the intended parameter is $\mu_d$, mean branch time saving with $d=\text{before}-\text{after}$. Measurements within a branch are matched, so an independent-samples test is inappropriate. There are nine observed differences, not 21 independent observations. Complete-case analysis may be biased because missingness is related to severe delays; it cannot automatically represent all 12 branches or a larger population. With nine pairs, an influential difference also threatens the approximate-normal-differences condition. Investigate data errors, explain missingness, recover measurements where possible, and inspect a difference plot before selecting an analysis with suitable support. Do not discard an extreme branch merely to obtain significance. No paired t decision is justified from this description alone.
\secInicio[prob:conditions][sec:advanced]
\end{document}

